% !TEX program = xelatex
\documentclass[11pt]{ctexart}

\usepackage[a4paper,top=2.6cm,bottom=2.6cm,left=2.8cm,right=2.8cm]{geometry}
\usepackage{amsmath,amssymb,amsthm}
\usepackage{booktabs}
\usepackage{array}
\usepackage{enumitem}
\usepackage{tikz}
\usetikzlibrary{arrows.meta,positioning}
\usepackage[colorlinks=true,linkcolor=blue!70!black,citecolor=blue!70!black,urlcolor=blue!70!black]{hyperref}
\usepackage[numbers,sort&compress]{natbib}

\theoremstyle{plain}
\newtheorem{theorem}{定理}[section]
\newtheorem{definition}[theorem]{定义}
\newtheorem{example}[theorem]{例}
\theoremstyle{remark}
\newtheorem{remark}[theorem]{说明}

\newcommand{\benchmark}{\textsc{Benchmark}}
\newcommand{\solution}{\textsc{Solution}}

\title{\bfseries 用 AI Agent 自动构造 Type Denotation / Logical Relation\\[6pt]
\large 研究计划：\benchmark{} 构建（与南京大学张子豪合作）与 \solution{} 探索}
\author{周喆（Zhe Zhou）}
\date{2026 年 8 月}

\begin{document}

\maketitle

\begin{abstract}
\noindent
本项目研究如何让 AI Agent 自动构造 PL / 程序验证论文中最核心的语义不变量
（type denotation、term denotation、logical relation、semantic domain），
把 AI for Proof 从``写 tactic''推进到``发现证明的核心不变量''。
本文档把工作拆成两部分：
\textbf{\benchmark{}}（近期，由南京大学张子豪同学负责）：构建``给定语言与定理、构造语义不变量''的
任务集，从类型系统推广到 abstract machine / definitional interpreter / staged evaluation 等更广义的 PL 语义工作；
\textbf{\solution{}}（后续，周喆负责）：agent 系统原型、证明失败反馈回路、
以及``哪些复杂性来自问题本身、哪些来自证明框架''的比较研究。
未来该方向天然连接 logical relation、denotational semantics、substructural logic 与 separation logic，
可作为与北京大学数学科学学院合作的切入点。
\end{abstract}

\tableofcontents

\section{核心想法与定位}
\label{sec:idea}

本项目（暂名 ``AI Denotation Agent''）的目标是：\emph{让 AI Agent 自动构造 PL / 程序验证论文中最关键的
语义不变量}，即 type denotation、term denotation、logical relation，以及在必要时选择合适的 semantic domain。

目标不是让 LLM 简单地把 proof script 补完，而是让它参与\emph{``为什么这个类型系统是 sound 的''}这一最核心的
\emph{设计}层面：在给定语言语法、操作语义、typing rules 与目标 soundness theorem 之后，自动提出候选的
denotation / logical relation，并迭代地对每条 typing rule 尝试语义 soundness 证明，在失败时分析原因
（不变量太强、太弱、缺少 resource、缺少 monotonicity、semantic domain 不够）并修改设计。

\paragraph{关键推广：不局限于类型系统。}
``构造正确的不变量''这一核心任务并不只出现在类型系统里。抽象机（abstract machine）、定义式解释器
（definitional interpreter）、程序变换（例如 staged evaluation 中的 let-insertion）等一大类 PL 语义工作，
其正确性（soundness、语义等价、语义保持）同样依赖一个精心选择的语义结构，而且这些工作大多只有非形式化论证，
甚至完全没有机械化形式化（详见 \S\ref{sec:benchmark}）。
因此本项目把 \benchmark{} 从``类型系统''推广到``一般 PL 语义工作''：既扩大了数据规模，
也使``语义不变量构造''成为一个真正通用的 AI for Proof 任务。

\paragraph{近期分工。}
\begin{itemize}
  \item \textbf{张子豪（南京大学）}：负责 \benchmark{} 构建与基线评估（\S\ref{sec:benchmark}）。
        这部分产出明确、可以立即启动：编写问题卡片、运行现有 LLM agent 作为基线、记录并整理结果。
  \item \textbf{周喆}：负责整体设计与 \solution{} 探索（\S\ref{sec:solution}）：
        agent 系统原型、失败反馈回路、证明框架负担的比较研究。
\end{itemize}

\paragraph{远期合作。}
该项目天然连接 logical relation、denotational semantics、substructural logic、separation logic 等
数学与逻辑内容，未来希望邀请\emph{北京大学数学科学学院}的同事加入（详见 \S\ref{sec:collab}），
形成``数院提供数学语义基础、系统 / PL 方向提供语言与验证问题、AI agent 提供自动化''的合作闭环。

\section{背景与动机}
\label{sec:background}

\subsection{为什么 type denotation / logical relation 是核心难点}

在 PL 与程序验证论文中，有一大类工作是在设计类型系统（例如 Rust 的 ownership / borrowing 保证内存安全）。
对简单语言，type soundness 通常可以用 progress + preservation 证明
\citep{WrightFelleisen1994}；但对更复杂的类型系统——Rust 的 lifetime、unsafe 抽象、并发、
substructural resource 等——纯语法证明往往不够，通常需要 denotational semantics 或 logical relation。

这类证明的核心难点，是\emph{如何定义合适的 type denotation / logical relation}。
它本质上就是程序验证中的 inductive invariant：
\begin{itemize}
  \item 定义太弱，推不出安全性；
  \item 定义太强，typing rule 无法证明 sound；
  \item 定义方式不合适，证明会非常复杂，甚至需要引入新的 semantic domain
        （separation logic、resource algebra、lifetime logic、step-indexed Kripke logical relation 等）。
\end{itemize}

两篇背景文献：
\begin{itemize}
  \item \citet{Timany2024JACM}：说明为什么 semantic / logical type soundness 比 progress + preservation
        更强，特别适合解释 abstraction 与 unsafe / safe 交互；
  \item \citet{Reynolds2000Meaning}：提供 intrinsic / extrinsic semantics 的经典视角——
        把类型看成程序语义上的 property / relation，而非只看 typing derivation 本身。
\end{itemize}

\subsection{与 AI for Proof 现状的差距}

目前很多 AI for Proof 工作集中在：证明单个 lemma、生成 tactic、proof repair、theorem proving benchmark。
但在 PL / 程序验证里，真正困难的地方经常不是某一步 tactic，而是：

\begin{quote}
\emph{我们到底应该定义什么样的 invariant，才能让整个 soundness proof 成立？}
\end{quote}

这个问题目前基本只能靠人的洞察力解决。如果能被 agent 部分自动化，将直接降低 PL mechanization 的门槛，
并让研究者更快地验证一个新的 typing rule 是否 sound、探索 alternative denotation。

\subsection{一个被忽视的事实：大量 PL 语义工作没有机械化形式化}

类型系统之外，还有一大类 PL 语义工作——abstract machine、definitional interpreter、staged evaluator 等——
它们论文里的正确性声明（soundness、语义等价、语义保持）通常只有非形式化论证，没有公开的机械化形式化
（详见 \S\ref{sec:machine-line}）。这些工作恰好是``不变量构造''任务最干净、也最有补全价值的数据来源。

\section{\benchmark{}：``语义不变量构造''任务集（张子豪）}
\label{sec:benchmark}

这一部分是可以立即启动的工作，由张子豪同学负责执行，周喆提供设计与工具链支持。
\benchmark{} 的目标是回答一个定量问题：\emph{给定语言定义、语义与目标定理，现有的 LLM agent
能在多大程度上自动构造出正确、可证明的语义不变量？}

\subsection{统一任务格式（问题卡片）}

每个任务用一张\emph{问题卡片}（problem card）描述，尽量机器可读、可复用：

\begin{quote}
\noindent\textbf{输入：}
\begin{itemize}
  \item 语言语法（BNF，或已有的 Coq / Lean 定义，优先复用现有 artifact）；
  \item 操作语义（big-step / small-step / abstract machine / definitional interpreter）；
  \item typing rules，或待验证的目标语义对象；
  \item 目标定理（type soundness / 抽象解释 soundness / 变换语义保持 / 机器与解释器等价）；
  \item 可选：例子、测试、已有 proof attempt。
\end{itemize}
\noindent\textbf{期望输出：}
\begin{itemize}
  \item semantic domain（或 abstract domain + Galois connection / abstraction relation）；
  \item type denotation / term denotation / logical relation（或 simulation relation）；
  \item context interpretation；
  \item 每条 typing rule（或每个语义步骤）对应的 proof obligation；
  \item 完整证明脚本（Rocq / Lean），要求无 \texttt{Admitted}。
\end{itemize}
\end{quote}

\subsection{三类任务（类型系统层）}

类型系统层的任务分为三类，分别对应不同研究目的：

\begin{enumerate}
  \item \textbf{复现与简化}：论文本身已有 denotation / logical relation。
        例如 RustBelt \citep{Jung2018RustBelt}、RustHornBelt \citep{Matsushita2022RustHornBelt}、
        VerusBelt \citep{Hance2026VerusBelt}，以及 reachability types 的逻辑关系工作
        \citep{Bao2025ReachabilityLR}——这些工作的 type denotation / logical relation 本身就是核心贡献。
        这里的目标不是假装从零发明，而是看 agent 能否根据 typing rules 与 soundness theorem 找到类似
        的不变量，或者发现更简单、更模块化的定义方式。
  \item \textbf{证明结构迁移}：论文有 soundness proof，但不是 denotation / logical relation
        （多为 progress + preservation，如 STLC、System F、Featherweight Java \citep{Igarashi2001FJ}）。
        目标：验证 agent 能否把 syntactic proof 问题重铸成 semantic type soundness / logical relation
        问题，而不是只会复述原论文。
  \item \textbf{补全缺失证明}：系统没有完整证明，或只有非形式化论证。输入只有语言定义、typing rules、
        例子与目标定理。成功则补全 soundness proof；失败也可能暴露 typing rule 本身不 sound，
        或指出缺少必要的语义结构。这类任务最接近真实研究场景。
\end{enumerate}

\subsection{推广：超越类型系统（abstract machine 线）}
\label{sec:machine-line}

原想法只覆盖``类型系统''；本项目的关键推广是：\emph{把``构造语义不变量''任务从类型系统推广到一般的
PL 语义工作}。理由有二：(1) 任务结构同构——都是``给定形式定义与目标定理，构造让证明成立的语义结构''；
(2) 数据更丰富——很多这类工作没有机械化形式化，benchmark 直接产生研究价值。

\begin{table}[htbp]
\centering
\caption{abstract machine 线候选任务（正确性声明多为非形式化论证，未见公开机械化形式化）}
\label{tab:machine-line}
\begin{tabular}{@{}>{\raggedright\arraybackslash}p{0.35\textwidth}>{\raggedright\arraybackslash}p{0.36\textwidth}>{\raggedright\arraybackslash}p{0.24\textwidth}@{}}
\toprule
\textbf{论文} & \textbf{正确性声明} & \textbf{要构造的语义结构} \\
\midrule
Van Horn \& Might, \emph{Abstracting Abstract Machines} (ICFP'10) \citep{VanHornMight2010AAM}
  & 抽象机上的抽象解释（CFA）相对于具体机器语义是 sound 的
  & abstract domain + abstraction / simulation relation \\
Johnson \& Van Horn, \emph{Abstracting Definitional Interpreters} (ICFP'14) \citep{JohnsonVanHorn2014ADI}
  & 把抽象机统一为 definitional interpreter 后，抽象版本仍 sound
  & monadic abstract interpreter 的 soundness invariant \\
Danvy \& Nielsen, \emph{Refocusing in Reduction Semantics} (RS-04-26, 2004) \citep{DanvyNielsen2004Refocusing}
  & refocus 与 plug-then-decompose 计算相同结果
  & 约简语义与抽象机之间的等价 / simulation \\
Danvy \& Millikin, \emph{Refunctionalization at Work} (SCP'09) \citep{DanvyMillikin2009Refunctionalization}
  & defunctionalization 的左逆在语义上保持计算
  & 函数表示与一等函数间的语义等价 \\
Wei, Decker \& Rompf, \emph{Refunctionalization of Abstract Abstract Machines} (ICFP'18, pearl) \citep{Wei2018RefunctionalizationAAM}
  & AAM 与 ADI 两个框架通过 refunctionalization 相互转换
  & 机器状态与解释器状态之间的双向转换保持 soundness \\
Wei, Tan \& Zhong, \emph{Let It Be Optimized: Multi-Stage Evaluators with Let-Insertion} (ICFP'26, pearl) \citep{Wei2026LetItBeOptimized}
  & staged evaluator 中 \mbox{let-insertion} 及各步优化均保持程序语义
  & 分阶段求值 / \mbox{let-insertion} 的语义保持不变量 \\
\bottomrule
\end{tabular}
\end{table}

\noindent
这些任务的``语义不变量''与类型系统的 logical relation 一一对应：
类型系统对应 type denotation / logical relation；机器与解释器对应 machine state 之间的（bi-）simulation；
变换正确性对应 refinement / equational theory。\citet{Wei2026LetItBeOptimized} 是该线最近的工作，
其 related work 展示了完整脉络（MetaML \citep{TahaSheard2000MetaML}、LMS \citep{RompfOdersky2010LMS}、
staged abstract interpreters \citep{Wei2019StagedAI} 等），可以作为问题卡片书写的文献地图。

\begin{remark}
  ``没有机械化形式化''这一判断需要在构建问题卡片时逐篇核实（搜索 artifact、公开 Coq/Lean 代码）。
  如果某篇已有形式化，就把它降级为第一类（复现与简化）任务——这本身也是 benchmark 元数据的一部分。
\end{remark}

\subsection{任务优先级与难度分级}

\begin{table}[htbp]
\centering
\caption{benchmark 优先级（P0 最先启动）}
\label{tab:priority}
\begin{tabular}{@{}p{0.06\textwidth}p{0.30\textwidth}p{0.36\textwidth}p{0.18\textwidth}@{}}
\toprule
\textbf{优先级} & \textbf{任务} & \textbf{说明} & \textbf{类别} \\
\midrule
P0 & STLC logical relation（call-by-value） & 校准任务：有大量现成 Coq/Lean 参考，用于验证 agent 管线与评价指标 & 一 \\
P0 & System F parametricity & 校准任务：逻辑等价 + 自由定理，验证 agent 能否处理高阶量词 & 一 \\
P1 & RustBelt 最小片段（$\lambda_{\mathrm{rust}}$ 子集，\texttt{own}/\texttt{shr} 模型） & 已有 Iris + Coq 形式化，测试 agent 能否恢复 \texttt{own}/\texttt{shr} split、current/future model & 一 \\
P1 & Reachability Types 片段 & OOPSLA'25 已给出 logical relation \citep{Bao2025ReachabilityLR}，POPL'24 有 Coq 机械化 \citep{Bao2024PolyReachability}；测试简化 / 重组织 & 一 \\
P2 & Featherweight Java 重铸 & 把 syntactic type soundness 重铸为 semantic soundness \citep{Igarashi2001FJ} & 二 \\
P2 & region / ownership 小语言 & 测试 agent 能否发现资源敏感（resource-sensitive）语义 & 二 / 三 \\
P3 & AAM / ADI 的 soundness 形式化 & 补全缺失证明，见 \tablename~\ref{tab:machine-line} 前两行 & 三 \\
P3 & Refunctionalization AAM $\leftrightarrow$ ADI 等价性 & 补全缺失证明，见 \tablename~\ref{tab:machine-line} 第五行 & 三 \\
P3 & let-insertion 语义保持（ICFP'26 pearl） & 最新、难度最高，见 \tablename~\ref{tab:machine-line} 末行 & 三 \\
\bottomrule
\end{tabular}
\end{table}

\subsection{交付物与里程碑}

每个 case 的交付物（\textbf{四件套}）：
\begin{enumerate}
  \item \textbf{问题卡片}：输入 / 输出 / 目标定理，机器可读（Markdown 或 JSON），附语言定义骨架（Rocq / Lean）；
  \item \textbf{运行日志}：agent 的 candidate denotation、失败尝试与原因分析、最终结果（成功的定义与证明）；
  \item \textbf{形式化产物}：Rocq / Lean 证明脚本（无 \texttt{Admitted}）+ 说明文档；
  \item \textbf{元数据}：难度、是否已有公开形式化、参考文献、备注（例如发现原规则过强/过弱）。
\end{enumerate}

\begin{table}[htbp]
\centering
\caption{建议里程碑}
\label{tab:milestones}
\begin{tabular}{@{}p{0.16\textwidth}p{0.30\textwidth}p{0.44\textwidth}@{}}
\toprule
\textbf{时间} & \textbf{里程碑} & \textbf{内容} \\
\midrule
第 1--2 周 & M0：跑通管线 & 完成 P0 两个校准 case 的问题卡片；用现有 LLM agent（如 Codex）+ Rocq 跑出基线；
             确定运行记录模板与评价指标 \\
第 3--6 周 & M1：类型系统层 & 完成 P1、P2 共 4--6 个 case；整理``复现 / 迁移 / 补全''三类结果对比 \\
第 7--12 周 & M2：推广层 & 完成 abstract machine 线 2--3 个 case（建议先 AAM，再 refunctionalization，
             最后 let-insertion）；产出中期报告 \\
持续 & 元数据维护 & 每篇候选论文核实是否有公开形式化，维护 benchmark 清单 \\
\bottomrule
\end{tabular}
\end{table}

\subsection{工具链建议}

\begin{itemize}
  \item \textbf{证明助手}：Rocq 9.1 + stdpp（团队已有成熟经验，见 \texttt{UnderLogicAndType} 项目）；
        Lean 4 作为对照后端，方便与``证明自动化''生态对接；
  \item \textbf{框架}：Iris / separation logic 仅在对齐原论文时使用（P1 的 RustBelt 片段）；
        其他 case 鼓励``直接定义''的轻量路线，以便比较框架负担（\S\ref{sec:framework}）；
  \item \textbf{自动化}：coqhammer、\texttt{lia}、\texttt{eauto}、QuickChick（反例搜索）等；
  \item \textbf{Agent 基线}：Codex / Claude 等现成 agent，统一 prompt 模板，记录每次运行。
\end{itemize}

\section{\solution{}：Agent 系统与科学问题（后续探索）}
\label{sec:solution}

\benchmark{} 为 \solution{} 提供数据与评价基准；\solution{} 反过来决定 \benchmark{} 的问题卡片该怎么出。
两者并行推进：\benchmark{} 先跑基线，\solution{} 在此之上迭代 agent 能力。

\subsection{系统工作流}

\begin{figure}[htbp]
\centering
\begin{tikzpicture}[
  node distance=0.9cm and 1.1cm,
  box/.style={draw, rounded corners=2pt, align=center, minimum width=3.4cm, font=\small},
  arr/.style={-{Stealth[length=2.5mm]}, thick},
]
  \node[box] (in) {输入：语言定义\\ + typing rules + 目标定理};
  \node[box, right=of in] (gen) {Agent 提议 candidate\\ semantic domain / denotation};
  \node[box, right=of gen] (prove) {尝试证明每条 rule\\ 的 semantic soundness};
  \node[box, below=of prove] (succ) {全部规则可证};
  \node[box, below=of gen] (fail) {失败：分析原因\\ 太强 / 太弱 / 缺 resource\\ 缺 monotonicity / domain 不够};
  \node[box, below=of in] (gen2) {修改 denotation 或 domain};

  \draw[arr] (in) -- (gen);
  \draw[arr] (gen) -- (prove);
  \draw[arr] (prove) -- node[right] {成功} (succ);
  \draw[arr] (prove) -- node[right] {失败} (fail);
  \draw[arr] (fail) -- (gen2);
  \draw[arr] (gen2) -- node[above] {迭代} (in);
  \draw[arr] (succ) -| node[above right, font=\small] {推广检查} (gen);
\end{tikzpicture}
\caption{Agent 迭代工作流：构造 denotation 的过程类似程序验证中寻找 inductive invariant。}
\label{fig:workflow}
\end{figure}

\noindent
收敛之后，进一步检查：\emph{某个 denotation 成功证明了原 typing rules，是否还允许更强、更自由的
typing rules？}如果是，说明原论文的规则可能过于保守。这个方向不仅能复现已有证明，
还可能发现已有类型系统设计中的限制或改进空间。

\subsection{关键技术挑战}

\begin{itemize}
  \item \textbf{候选生成}：如何让 LLM 生成``有前途''的 denotation？
        把 denotation 的常见结构（product / sum / function / existential / step-indexing /
        resource algebra / Kripke world）做成一个小的组合 DSL，让 agent 在此空间内搜索，
        而不是自由文本生成；
  \item \textbf{失败反馈回路}：把证明失败（未闭合目标、循环依赖、不可判定的子目标）转成可操作的信号；
        用 QuickChick / small-model 搜索找反例，区分``denotation 错了''与``证明还差一步''；
  \item \textbf{证明后端}：Rocq / Lean 双后端；证明自动化（coqhammer、\texttt{lia}、\texttt{eauto}）
        与 denotation 生成的解耦，使失败分析聚焦于\emph{不变量设计}而非 tactic 细节；
  \item \textbf{评价指标}：成功率之外，还要度量 denotation 的\emph{可读性 / 模块化}、
        是否需要引入新的 semantic domain、能否发现更强的 typing rules。
\end{itemize}

\subsection{科学问题：证明框架负担与 denotation 发现}
\label{sec:framework}

很多已有证明依赖大型框架（如 Iris \citep{Jung2018Iris}、separation logic 库、step-indexed LR 库）。
这些框架表达力强，但也可能带来额外负担：某些 denotation 必须绕着框架接口写，某些概念要通过复杂的
resource algebra 或 proof pattern 间接表达，导致定义不直观、proof script 很重、可读性差。

如果 agent 能在同一个 soundness 任务上尝试不同的 semantic domain、不同的 denotation factoring、
不同的 proof architecture，我们就可以系统地比较：
\begin{itemize}
  \item 哪些复杂性来自问题本身，哪些来自当前证明框架的限制；
  \item 是否存在更直接、更可读的 denotation；
  \item 是否需要为这类证明设计新的 library abstraction 或新的 semantic framework。
\end{itemize}
因此这个项目不只是评估 agent 能否证明定理，也可以反过来帮助评估 Iris 等框架在表达某类
type denotation / logical relation 时的工程负担，并为设计更好的 proof harness 提供指导。

\subsection{为什么这个问题重要}

\begin{enumerate}
  \item \textbf{帮助设计复杂类型系统}：研究者可以更快验证 typing rule 是否 sound，探索 alternative denotation；
  \item \textbf{降低 PL mechanization 门槛}：很多论文形式化证明的难点在 denotation / logical relation 设计；
  \item \textbf{发现已有系统的不足}：agent 可能发现原 denotation 过于复杂、原 typing rules 过于保守，
        甚至指出某些规则实际不 sound；
  \item \textbf{连接数学逻辑与系统验证}：天然结合 logical relation、denotational semantics、
        substructural logic、separation logic 与程序验证；
  \item \textbf{RustBelt 风格 case study 有影响力}：若能在 RustBelt 风格 fragment 上展示 agent 自动发现
        或简化 denotation，会很有说服力，也容易获得 PL 社区关注；
  \item \textbf{改进现有证明框架}：见 \S\ref{sec:framework}。
\end{enumerate}

\section{路线图与分工}
\label{sec:roadmap}

\begin{table}[htbp]
\centering
\caption{阶段划分与负责人}
\label{tab:roadmap}
\begin{tabular}{@{}p{0.14\textwidth}p{0.30\textwidth}p{0.30\textwidth}p{0.16\textwidth}@{}}
\toprule
\textbf{阶段} & \textbf{时间} & \textbf{内容} & \textbf{负责人} \\
\midrule
阶段 0 & 2026-08 起 & \benchmark{} 构建：问题卡片、基线运行、元数据（\S\ref{sec:benchmark}） & 张子豪（周喆指导） \\
阶段 1 & 2026 秋 & \solution{} 原型 v1：针对 P0/P1 的 denotation 候选生成 + 迭代证明 & 周喆 \\
阶段 2 & 2027 春 & 推广到 abstract machine 线；框架负担比较实验 & 张子豪 + 周喆 \\
阶段 3 & 2027 年 & 论文产出：benchmark 报告 / 方法论文；评估与北大数院合作 & 全体 \\
\bottomrule
\end{tabular}
\end{table}

\subsection{与北京大学数学科学学院的合作}
\label{sec:collab}

这个方向为数学侧提供了明确的问题来源：logical relation 的模型论基础、denotational semantics 中的
domain theory、substructural logic 与 categorical semantics、step-indexing 与 Kripke 语义。
反过来，数院可以为 denotation 的``正确形式''提供理论指导，例如：
\begin{itemize}
  \item denotation 的范畴 / 类型论表述（intrinsic vs.\ extrinsic semantics \citep{Reynolds2000Meaning}）；
  \item 不同 resource algebra 之间能否建立系统的比较与合成；
  \item ``最优 / 最简 denotation''能否在某个数学意义下被刻画（例如作为某个 universal property）。
\end{itemize}
这些问题的答案，既能让 agent 的搜索更有方向，也能把``AI 自动构造不变量''做成一个有数学深度的研究方向。

\subsection{立即可以开始的事情}

\begin{enumerate}
  \item 张子豪按 \S\ref{sec:benchmark} 完成 P0 两张问题卡片，并跑通一次 Codex/Rocq 基线；
  \item 维护候选论文清单（类型层 + machine 层），逐篇核实是否有公开形式化；
  \item 固定运行记录模板与评价指标，保证后续结果可复现、可比较。
\end{enumerate}


\bibliographystyle{plainnat}
\bibliography{bibliography}

\end{document}
